\section{Mathematical Derivation of Variational Free Energy}
\label{appendix:derivation}

In the context of Active-Bidder, we model the bidding agent as a self-evidencing system that minimizes its Variational Free Energy (VFE) to maintain homeostasis within a volatile auction environment. Here, we provide the step-by-step derivation of the VFE from the fundamental definition of surprise.

\subsection{Surprise and the Evidence Lower Bound (ELBO)}
Let $\mathbf{o} \in \mathcal{O}$ denote the sequence of observations (e.g., auction outcomes, click-through signals) and $\mathbf{s} \in \mathcal{S}$ denote the latent states of the environment (e.g., hidden user intent, competitor strategies). The goal of the agent is to minimize the ``surprise'' or negative log-evidence:
\begin{equation}
    -\ln p(\mathbf{o}) = -\ln \int p(\mathbf{o}, \mathbf{s}) d\mathbf{s}
\end{equation}
Directly minimizing surprise is computationally intractable due to the integral over the high-dimensional latent space $\mathcal{S}$. We introduce a recognition density $q(\mathbf{s})$, which represents the agent's internal belief about the hidden states. Using Jensen's Inequality, we derive the Evidence Lower Bound (ELBO):
\begin{align}
    -\ln p(\mathbf{o}) &= -\ln \int q(\mathbf{s}) \frac{p(\mathbf{o}, \mathbf{s})}{q(\mathbf{s})} d\mathbf{s} \\
    &\leq -\int q(\mathbf{s}) \ln \frac{p(\mathbf{o}, \mathbf{s})}{q(\mathbf{s})} d\mathbf{s} \triangleq \mathcal{F}(q, \mathbf{o})
\end{align}
where $\mathcal{F}$ is the Variational Free Energy. By definition, $\mathcal{F} \geq -\ln p(\mathbf{o})$, meaning minimizing $\mathcal{F}$ provides an upper bound on surprise.

\subsection{Functional Decompositions of VFE}
The VFE can be decomposed into two insightful forms that guide the bidding strategy.

\textbf{1. Divergence + Energy Formulation:}
We can rewrite the VFE by separating the recognition density from the posterior:
\begin{align}
    \mathcal{F} &= \int q(\mathbf{s}) \ln \frac{q(\mathbf{s})}{p(\mathbf{o}, \mathbf{s})} d\mathbf{s} \\
    &= \int q(\mathbf{s}) \ln \frac{q(\mathbf{s})}{p(\mathbf{s}|\mathbf{o})p(\mathbf{o})} d\mathbf{s} \\
    &= D_{KL}[q(\mathbf{s}) \| p(\mathbf{s}|\mathbf{o})] - \ln p(\mathbf{o})
\end{align}
This shows that when the VFE is minimized, the recognition density $q(\mathbf{s})$ becomes a perfect approximation of the true posterior $p(\mathbf{s}|\mathbf{o})$, and the VFE itself converges to the surprise.

\textbf{2. Complexity + Accuracy Formulation:}
Alternatively, we can decompose $\mathcal{F}$ into terms representing the cost of updating beliefs and the fit to observations:
\begin{align}
    \mathcal{F} &= \int q(\mathbf{s}) \ln \frac{q(\mathbf{s})}{p(\mathbf{o}|\mathbf{s})p(\mathbf{s})} d\mathbf{s} \\
    &= \int q(\mathbf{s}) \ln \frac{q(\mathbf{s})}{p(\mathbf{s})} d\mathbf{s} - \int q(\mathbf{s}) \ln p(\mathbf{o}|\mathbf{s}) d\mathbf{s} \\
    &= \underbrace{D_{KL}[q(\mathbf{s}) \| p(\mathbf{s})]}_{\text{Complexity}} - \underbrace{\mathbb{E}_{q(\mathbf{s})}[\ln p(\mathbf{o}|\mathbf{s})]}_{\text{Accuracy}}
\end{align}
In the bidding context, \textit{Complexity} penalizes beliefs that deviate excessively from the prior (preventing overfitting to transient auction noise), while \textit{Accuracy} ensures the agent's model correctly predicts market outcomes.

\subsection{Expected Free Energy (G) for Policy Selection}
For future-oriented decision making (bidding), the agent selects a policy $\pi$ (a sequence of bids) that minimizes the Expected Free Energy (EFE), denoted as $G(\pi)$. Unlike VFE, EFE is calculated over future observations $\mathbf{o}_\tau$ that have not yet occurred:
\begin{equation}
    G(\pi) \approx \sum_\tau \mathbb{E}_{q(\mathbf{o}_\tau, \mathbf{s}_\tau | \pi)} [\ln q(\mathbf{s}_\tau | \pi) - \ln p(\mathbf{o}_\tau, \mathbf{s}_\tau)]
\end{equation}
This can be further expanded into risk and ambiguity terms:
\begin{equation}
    G(\pi) \approx \underbrace{D_{KL}[q(\mathbf{o}_\tau | \pi) \| p(\mathbf{o}_\tau)]}_{\text{Risk (Relative Entropy)}} + \underbrace{\mathbb{E}_{q(\mathbf{s}_\tau | \pi)}[H(p(\mathbf{o}_\tau | \mathbf{s}_\tau))]}_{\text{Ambiguity}}
\end{equation}
where $p(\mathbf{o}_\tau)$ represents the agent's ``prior preferences'' (e.g., achieving a target ROI). Risk-minimization drives the agent toward preferred outcomes, while ambiguity-minimization drives exploratory bidding to resolve uncertainty about the market state.

\section{Implementation Details and Reproducibility}
\label{appendix:implementation}

The Active-Bidder framework is implemented using PyTorch 2.1.2 and the \texttt{Pyro} probabilistic programming library for variational inference.

\subsection{Infrastructure and Hardware}
All experiments were conducted on a high-performance computing cluster. The primary training specifications are:
\begin{itemize}
    \item \textbf{GPU:} $8 \times$ NVIDIA A100 (80GB SXM4) interconnected via NVLink.
    \item \textbf{CPU:} Dual AMD EPYC 7742 64-Core Processors.
    \item \textbf{RAM:} 512GB DDR4.
    \item \textbf{Software:} CUDA 12.1, Python 3.10, PyTorch Lightning for distributed training.
\end{itemize}

\subsection{Hyperparameter Configuration}
Table \ref{tab:hyperparams} lists the specific hyperparameters used for the results reported in Section 5. We utilized the AdamW optimizer \cite{loshchilov2017decoupled} with a cosine annealing learning rate schedule.

\begin{table}[h]
\centering
\caption{Hyperparameter settings for Active-Bidder training.}
\label{tab:hyperparams}
\resizebox{\columnwidth}{!}{%
\begin{tabular}{l|l|l}
\hline
\textbf{Category} & \textbf{Parameter} & \textbf{Value} \\ \hline
Optimization & Learning Rate ($\eta$) & $3 \times 10^{-4}$ \\
& Weight Decay & $1 \times 10^{-2}$ \\
& Batch Size & 4096 \\
& Optimizer & AdamW ($\beta_1=0.9, \beta_2=0.999$) \\ \hline
Model Architecture & Latent Dim ($|s|$) & 64 \\
& Hidden Units (MLP) & [512, 256, 128] \\
& Activation Function & Swish (SiLU) \\
& Dropout & 0.1 \\ \hline
FEP Specifics & Observation Noise ($\sigma$) & 0.05 (learnable) \\
& Prior ROI Preference & 4.5 \\
& Horizon ($T$) & 24 steps (hourly) \\
& $\beta$ (VFE Weight) & 1.0 (standard) \\ \hline
\end{tabular}
}%
\end{table}

\subsection{Feature Engineering}
The model ingests a multi-modal feature vector $\mathbf{x}_t$ consisting of historical auction logs and real-time context. The feature categorization is detailed in Table \ref{tab:features}.

\begin{table}[h]
\centering
\caption{Feature Space for the Generative Model.}
\label{tab:features}
\resizebox{\columnwidth}{!}{%
\begin{tabular}{lll}
\hline
\textbf{Feature Group} & \textbf{Description} & \textbf{Type} \\ \hline
\textit{User Context} & Browser, OS, Geolocation, Device Type & Categorical (Embedding) \\
\textit{Auction State} & Winning Price (t-1), Second Bid, Floor Price & Numerical (Normalized) \\
\textit{Temporal} & Hour of Day, Day of Week, Is Holiday & Cyclic Encoding \\
\textit{Performance} & Moving Average CTR (1h, 6h, 24h) & Numerical \\
\textit{Ad Attributes} & Creative ID, Campaign Category, Advertiser ID & Categorical (Embedding) \\ \hline
\end{tabular}
}%
\end{table}

\section{Additional Results: Robustness to $\beta$-Weighting}
\label{appendix:results}

In variational inference, a common modification is the $\beta$-VAE formulation, which scales the KL-divergence (Complexity) term:
\begin{equation}
    \mathcal{F}_\beta = \beta \cdot D_{KL}[q(\mathbf{s}) \| p(\mathbf{s})] - \mathbb{E}_{q(\mathbf{s})}[\ln p(\mathbf{o}|\mathbf{s})]
\end{equation}
We evaluated the sensitivity of Active-Bidder to $\beta \in [0.1, 5.0]$ to understand the trade-off between belief flexibility and observational accuracy.

\subsection{Performance Metrics under Variable $\beta$}
Figure \ref{fig:beta_sensitivity} (omitted here, summarized in text) illustrates that for $\beta < 1.0$, the agent becomes ``over-confident,'' leading to high variance in bidding behavior and frequent budget exhaustion due to aggressive exploration. Conversely, for $\beta > 2.0$, the agent becomes ``over-cautious'' (prior-dominated), failing to adapt to sudden shifts in market floor prices.

\begin{itemize}
    \item \textbf{Low $\beta$ (0.1 - 0.5):} High Return on Ad Spend (ROAS) in stable environments, but 40\% higher probability of catastrophic budget depletion in volatile sessions.
    \item \textbf{Optimal $\beta$ (0.8 - 1.2):} The ``Homeostatic Zone'' where the agent maintains a steady ROI of $4.2 \pm 0.15$ while effectively managing liquidity across the 24-hour horizon.
    \item \textbf{High $\beta$ (> 2.0):} Significant reduction in win rate (down 25\%) as the agent's beliefs are too rigid to follow the shifting market distribution.
\end{itemize}

\subsection{Stability Analysis}
We measured stability using the Coefficient of Variation (CV) of the bid price. Active-Bidder with $\beta=1.0$ achieved a CV of 0.22, compared to 0.45 for a standard Reinforcement Learning (PPO) baseline. This confirms that the FEP-based objective naturally induces a ``smoothness'' in bidding strategies, which is highly desirable for real-world demand-side platforms (DSPs) to maintain predictable spending patterns.

\subsection{Ablation on Recognition Density}
We further tested the impact of the recognition density's functional form. Using a Multivariate Normal distribution with a diagonal covariance matrix provided the best balance between computational speed and expressive power. Attempts to use Normalizing Flows (NFs) for $q(\mathbf{s})$ yielded marginal gains in accuracy ($< 1\%$) while increasing inference latency by $4\times$, rendering them unsuitable for the 10ms response time required in Real-Time Bidding (RTB).